\chapter{Folate Test}

\section*{Definition and Overview}
The folate test is a blood test that measures the level of folate (vitamin B9) in the body, used to diagnose folate deficiency and to help find the cause of anaemia. There are two main measurements: serum folate, which reflects recent intake, and red cell folate, which reflects longer-term stores. The test is most often ordered alongside a full blood count and vitamin B12 measurement when a person has fatigue, anaemia, or a sore tongue, and it is part of the work-up of the large red cells (macrocytosis) seen in megaloblastic anaemia. Folate testing is also used to monitor people at risk of deficiency, including those with malabsorption and those on interfering medicines. This chapter explains when the test is used, what it involves, and how results are interpreted.

\section*{Epidemiology}
Folate testing is commonly performed in Australia, mainly in the investigation of anaemia and fatigue. Because bread flour is fortified with folic acid, population folate levels are generally adequate, and the proportion of blood tests showing true deficiency is relatively low. The test is most useful in specific groups: people with macrocytic anaemia, women planning pregnancy, people with coeliac disease or other malabsorption, heavy drinkers, and people taking medicines such as methotrexate. Overdiagnosis and unnecessary supplementation are more common problems than missed deficiency in the general population.

\section*{Aetiology and Pathophysiology}
Understanding the test requires understanding what the measurements mean.
\begin{itemize}
    \item \textbf{Serum folate:} measures folate in the blood at the time of testing; it falls within days of reduced intake, so it reflects recent dietary intake
    \item \textbf{Red cell folate:} measures folate within red blood cells, reflecting the average folate level over the preceding months
    \item \textbf{The link to anaemia:} folate is required for red cell production; deficiency causes megaloblastic anaemia with large red cells
    \item \textbf{The B12 link:} vitamin B12 deficiency causes an identical anaemia, and B12 is measured with folate so the two are distinguished
    \item \textbf{The masking problem:} giving folate to a person with B12 deficiency can correct the anaemia while the B12 nerve damage progresses, which is why B12 must be checked first
    \item \textbf{Factors affecting the result:} recent food intake raises serum folate, alcohol lowers it, and haemolysis in the sample can falsely raise it
    \item \textbf{Fasting:} some laboratories recommend testing after fasting, though this is not always required
\end{itemize}

\section*{Clinical Features}
The test is used when specific features are present.
\begin{itemize}
    \item Fatigue, weakness, and pallor suggesting anaemia
    \item A full blood count showing large red cells (macrocytosis)
    \item A sore, red tongue, mouth ulcers, or cracks at the corners of the mouth
    \item Risk factors: poor diet, alcohol excess, malabsorption, pregnancy, or folate-interfering medicines
    \item The test is a simple blood sample from a vein, requiring no special preparation beyond possibly fasting
    \item Results are usually available within days
    \item The test is often combined with B12 measurement and iron studies in the anaemia work-up
\end{itemize}

\section*{Differential Diagnosis}
A low or borderline folate result must be interpreted in context.
\begin{itemize}
    \item True folate deficiency from diet, malabsorption, or increased demand
    \item B12 deficiency, which needs separate testing and treatment
    \item A falsely low serum folate from recent alcohol intake or inadequate fasting
    \item A falsely normal result despite deficiency in some laboratory conditions
    \item Anaemia from iron deficiency, chronic disease, or other causes, which produces different blood features
    \item The clinical picture, not the number alone, determines treatment
\end{itemize}

\section*{When to Seek Medical Care}
Folate testing is arranged by a doctor, and the results should be discussed with the doctor who ordered them. Anyone with persistent fatigue, anaemia, or risk factors for deficiency should seek assessment rather than self-prescribing supplements.

\textbf{Call 000 (or local emergency services) for:}
\begin{itemize}
    \item Severe breathlessness, chest pain, or collapse from profound anaemia
    \item Sudden weakness or neurological symptoms
    \item Any emergency symptoms
\end{itemize}

\section*{Diagnosis and Investigations}
\begin{itemize}
    \item \textbf{Full blood count and film:} assesses for macrocytosis and other anaemia features
    \item \textbf{Serum folate:} the first-line test for recent status
    \item \textbf{Red cell folate:} ordered when long-term status is the question
    \item \textbf{Vitamin B12:} measured in the same work-up to exclude B12 deficiency
    \item \textbf{Additional tests:} iron studies, liver function, and thyroid function as indicated
    \item \textbf{Investigations for the cause:} coeliac serology and dietary review when malabsorption is suspected
\end{itemize}

\section*{Management}
\begin{itemize}
    \item \textbf{Normal results:} no treatment needed; dietary advice to maintain adequate intake
    \item \textbf{Low serum folate:} oral folic acid, typically 500 micrograms daily, with a repeat blood count to confirm recovery
    \item \textbf{Deficiency with anaemia:} treatment is followed by monitoring of the blood count over weeks
    \item \textbf{Pregnancy planning:} all women planning pregnancy take 500 micrograms of folic acid daily regardless of the test result
    \item \textbf{Treat the cause:} address diet, alcohol, malabsorption, or medicines
    \item \textbf{Monitor at-risk people:} repeat testing in those with ongoing risk factors
\end{itemize}

\section*{Complications}
\begin{itemize}
    \item Missed B12 deficiency if folate is treated without checking B12
    \item Unnecessary treatment and anxiety when borderline results are overinterpreted
    \item Deficiency continuing when the underlying cause is not treated
    \item Neural tube defects when deficiency in early pregnancy goes unrecognised
    \item Falsely reassuring results from improper sample handling
\end{itemize}

\section*{Prognosis and Outcomes}
The folate test is reliable when interpreted with the full blood count and B12, and treated deficiency resolves completely. Most abnormal results are corrected with a short course of folic acid and dietary change, and blood counts return to normal within weeks. For people at risk of deficiency, regular monitoring prevents recurrence. Because testing is combined with B12 measurement, the serious error of masking B12 deficiency is avoided. The test's value lies in accurate diagnosis, which leads to simple, effective treatment and excellent outcomes.

\section*{Prevention and Self-Care}
\begin{itemize}
    \item Discuss the purpose of the test with the doctor before having it
    \item Follow any fasting instructions from the laboratory
    \item Eat a diet rich in folate, and take folic acid when planning pregnancy
    \item Never take folate supplements to treat anaemia without medical advice
    \item Ensure B12 is checked at the same time
    \item Attend follow-up blood tests to confirm recovery
    \item Report new symptoms such as fatigue, sore tongue, or breathlessness
\end{itemize}

\section*{Sources and Further Reading}
The following reputable sources provide further information for healthcare students and patients seeking reliable, current advice about the folate test.
\begin{itemize}
    \item Healthdirect Australia. \textit{Folate test.} https://www.healthdirect.gov.au/
    \item Pathology Tests Explained. \textit{Folate.} https://pathologytestsexplained.org.au/
    \item The Royal College of Pathologists of Australasia. \textit{Folate testing.} https://www.rcpa.edu.au/
    \item NHS. \textit{Vitamin B12 or folate deficiency anaemia.} https://www.nhs.uk/conditions/vitamin-b12-or-folate-deficiency-anaemia/
    \item Mayo Clinic. \textit{Folate blood test.} https://www.mayoclinic.org/
    \item MedlinePlus. \textit{Folate test.} https://medlineplus.gov/
\end{itemize}
